\chapter{Résultats et discussion}

Dans ce chapitre, nous examinons la répartition des composantes d'\HI\ cinématiquement froide et chaude dans les galaxies proches, ainsi que leurs liens avec le gaz moléculaire et la formation stellaire.

\section[Répartition spatiale des composantes cinématiquement froide et chaude]{Répartition spatiale des composantes cinématiquement froide et chaude}
\label{sec:phase_spatial_distribution}

La figure~\ref{fig:individual_galaxies} présente les cartes de densité de surface de l'\HI\ total, cinématiquement froid et cinématiquement chaud des galaxies individuelles. Les deux composantes s'étendent sur une grande partie des disques atomiques cartographiés et présentent des condensations et des structures allongées au sein de l'émission étendue.

La composante cinématiquement froide représente en médiane environ 45\% de la masse d'\HI\ cartographiée, contre 55\% pour la composante chaude. À l'intérieur du disque optique ($r<r_{25}$), la composante chaude domine le bilan de masse d'\HI, même si la composante froide en constitue une part importante.

La composante froide présente un aspect plus filamentaire que la composante chaude : des structures fines et allongées se détachent de l'émission atomique plus diffuse. Cela conduit à se demander si le gaz froid se trouve préférentiellement dans les bras spiraux. Une telle préférence pourrait résulter de la compression du gaz dans les bras, qui augmente sa densité et favorise son refroidissement \citep{KimSpiral2008}. Pour repérer les bras des neuf galaxies de notre échantillon présentant une structure spirale visible, nous suivons la méthode de \citet{Foyle2010}, fondée sur l'émission à 3.6 $\mu$m pour tracer les bras stellaires. Nous comparons ensuite la fraction étroite dans les bras et les régions interbras. Cette fraction n'est pas systématiquement plus élevée dans les bras.

Les modèles d'équilibre thermique de l'\HI\ prévoient la coexistence d'un gaz froid et dense et d'un gaz tiède et diffus sur un intervalle de pressions thermiques \citep{Wolfire2003}. À pression comparable, le gaz froid est plus dense en raison de sa température plus basse. Les simulations montrent qu'une augmentation de densité du gaz atomique tiède peut permettre son refroidissement et sa condensation en structures froides \citep{Saury2014}. Le gaz froid peut alors fournir de la matière aux nuages moléculaires \citep{Krumholz2009Transition}. Le H$_2$ se forme principalement à la surface des grains de poussière, et sa formation est plus rapide aux fortes densités \citep{GloverMacLow2007}. Le rayonnement ultraviolet peut toutefois dissocier les molécules. Le gaz moléculaire s'accumule lorsque la poussière et le H$_2$ lui-même absorbent suffisamment de rayonnement pour le protéger. La transition atomique--moléculaire dépend ainsi de la densité du gaz, de la densité de colonne assurant cette protection et de l'abondance de poussière \citep{Krumholz2009Transition,Sternberg2014}. Dans les simulations de \citet{GloverMacLow2007}, le gaz initialement atomique se refroidit rapidement et la compression turbulente crée des régions denses où le H$_2$ se forme plus vite. Le lien entre gaz atomique froid et formation des nuages moléculaires repose donc sur les fortes densités qu'il peut atteindre, associées à une protection suffisante. Cela motive la recherche d'une association entre la composante étroite de l'\HI\ et l'émission du CO.

Nous comparons pour cela la fraction étroite dans des régions de même aire, nettement brillantes ou faibles en CO, à des rayons galactocentriques similaires, dans huit galaxies dont l'émission CO permet ce test. Nous partons des cartes natives d'\HI\ et de CO et les convoluons pour obtenir le même faisceau. Les résolutions spatiales résultantes vont de 234 à 721 pc.

Les deux composantes atomiques peuvent être renforcées près de l'émission CO brillante, mais la fraction froide n'y est pas systématiquement plus élevée. Le cas le plus net est NGC 5194 : la fraction froide moyenne pondérée par l'aire atteint environ 45\% dans les régions brillantes en CO, contre 29\% dans les régions voisines faibles en CO. Dans NGC 6946, les deux composantes augmentent près du CO brillant, tandis que la fraction froide reste voisine de 30\%. Les autres galaxies fournissent trop peu de régions appariées pour permettre une comparaison concluante. Nous ne trouvons donc pas d'indication systématique que l'\HI\ étroit soit davantage concentré près de l'émission CO que l'\HI\ large.

L'absence d'augmentation systématique de la fraction froide n'exclut pas que le gaz atomique froid alimente les nuages moléculaires. Dans les modèles de transition atomique--moléculaire, le gaz atomique demeure dans les couches qui entourent et protègent les intérieurs moléculaires, tandis que le gaz supplémentaire peut devenir moléculaire \citep{Krumholz2009Transition,Sternberg2014}. Une fois transformé en H$_2$, l'hydrogène ne contribue plus à l'émission d'\HI. Un réservoir moléculaire croissant ne s'accompagne donc pas nécessairement d'une contribution croissante de l'\HI\ étroit. De plus, le H$_2$ peut subsister là où le CO est dissocié \citep{Wolfire2010} : les régions de comparaison faibles en CO ne sont pas nécessairement dépourvues de H$_2$. Ces effets peuvent atténuer le contraste entre régions brillantes et faibles en CO. Notre comparaison ne permet donc pas d'établir si la composante étroite trace préférentiellement les régions où le gaz atomique devient moléculaire. Il reste difficile de localiser la transition \HI--H$_2$ dans notre échantillon, puisqu'une région peut contenir du H$_2$ sans émission CO détectable \citep{Wolfire2010}.

%\section[Méthode d'ajustement des relations de Kennicutt--Schmidt]{Méthode d'ajustement des relations de Kennicutt--Schmidt}
\label{sec:ks_fitting}
\label{sec:sflaw}

%Dans cette section, nous cherchons à déterminer laquelle des composantes d'\HI\ cinématiquement froide et chaude est la plus étroitement liée à la formation stellaire. Nous ajustons la relation de Kennicutt--Schmidt pour chaque composante gazeuse et comparons les pentes et les coefficients de corrélation, en prenant l'\HI\ total et le H$_2$ tracé par le CO comme références. Nous examinons les différences entre galaxies individuelles, ainsi qu'entre galaxies naines et spirales. La méthode d'ajustement est présentée dans la section~\ref{sec:ks_fitting_method}, puis les relations gaz--SFR dans la section~\ref{sec:ks_individual_results}.

%\subsection[Méthode d'ajustement]{Méthode d'ajustement}
\label{sec:ks_fitting_method}
Pour déterminer laquelle des composantes d'\HI\ cinématiquement froide et chaude est la plus étroitement liée à la formation stellaire, nous ajustons la relation de Kennicutt--Schmidt pour chaque composante gazeuse et comparons les pentes et les coefficients de corrélation, en prenant l'\HI\ total et le H$_2$ tracé par le CO comme références.
Suivant \citet{Fbigiel2008}, nous comparons les pentes des relations de Kennicutt--Schmidt des différentes composantes gazeuses. En coordonnées logarithmiques, la relation de l'équation (\ref{eq:SFLaw}) s'écrit :
\begin{equation}
\label{eq:ks_log_relation}
\log_{10}\Sigma_{\mathrm{SFR}} = a + N\log_{10}\Sigma_{\mathrm{gas}},
\end{equation}
\noindent où $a=\log_{10}A$ est l'ordonnée à l'origine et $N$ l'exposant de la loi de puissance.

\citet{Fbigiel2008} ont utilisé la bissectrice des deux régressions par moindres carrés ordinaires (OLS). Une régression OLS minimise les résidus verticaux ou horizontaux, selon la variable choisie comme variable dépendante. La méthode de la bissectrice combine ces deux régressions et traite les axes de façon symétrique. Or plusieurs distributions du gaz atomique de notre échantillon sont très dispersées ou presque verticales. Nous adoptons donc une régression orthogonale, qui minimise les distances perpendiculaires à la droite ajustée et estime directement l'orientation de la distribution.

Les pixels voisins à l'intérieur d'un faisceau contenant des mesures corrélées de la même émission, nous échantillonnons les cartes à des positions espacées d'au moins une largeur de faisceau, horizontalement et verticalement. Au préalable, les cartes de gaz et de SFR sont ramenées à la même résolution spatiale et à la même grille de pixels. À chaque position retenue, nous relevons les valeurs du gaz et du SFR dans le pixel correspondant. Nous ajustons la relation de Kennicutt--Schmidt à ces mesures, représentées par les points bleus de la figure~\ref{fig:individual_galaxies}. La distribution grisée montre l'ensemble des valeurs positives des pixels du champ cartographié.

Chaque mesure ainsi échantillonnée possède les coordonnées $(x_i,y_i)$, avec $x_i=\log_{10}\Sigma_{\mathrm{gas},i}$ et $y_i=\log_{10}\Sigma_{\mathrm{SFR},i}$, pour $i=1,\ldots,n$. Le nombre de mesures est noté $n$. Nous utilisons le critère de régression orthogonale \citep{Pearson1901,Isobe1990}, avec les poids définis ci-dessous. La droite ajustée minimise la somme pondérée des carrés des distances perpendiculaires :
\begin{equation}
\label{eq:ks_weighted_distance}
\mathcal{Q}(a,N)
=\sum_{i=1}^{n} w_i\,d_{\perp,i}^{\,2}
=\sum_{i=1}^{n} w_i\,\frac{(y_i-a-Nx_i)^2}{1+N^2},
\end{equation}
\noindent où $N$ est la pente de la relation de Kennicutt--Schmidt dans l'équation~\ref{eq:ks_log_relation}, $a$ l'ordonnée à l'origine et $w_i$ le poids attribué à la mesure $i$. La quantité $d_{\perp,i}$ est la plus courte distance entre le point $(x_i,y_i)$ et la droite $y=a+Nx$ dans le diagramme de $\log_{10}\Sigma_{\mathrm{SFR}}$ en fonction de $\log_{10}\Sigma_{\mathrm{gas}}$.

Afin de limiter l'influence des points dispersés dans les régions peu échantillonnées du diagramme gaz--SFR, nous attribuons davantage de poids aux points entourés de nombreuses mesures et moins aux points dont le voisinage en contient peu. La densité locale des points est estimée à l'aide d'un noyau gaussien bidimensionnel \citep{Scott1992}. La position de chaque point bleu $i$ dans le diagramme logarithmique est définie par :
\begin{equation}
\mathbf{z}_i =
\begin{pmatrix}
x_i \\
y_i
\end{pmatrix}.
\end{equation}
Son poids est proportionnel à la densité des autres points bleus dans son voisinage :
\begin{equation}
\label{eq:ks_density_weight}
w_i = \sum_{\substack{j=1\\j\ne i}}^{n}
\exp\!\left[-\frac{1}{2}
(\mathbf{z}_i-\mathbf{z}_j)^{\mathrm{T}}
\mathbf{H}^{-1}
(\mathbf{z}_i-\mathbf{z}_j)\right],
\end{equation}
\noindent où la somme porte sur tous les autres points bleus. La matrice $\mathbf{H}$ fixe l'étendue et la forme du voisinage gaussien. La règle de Scott \citep{Scott1992} donne $\mathbf{H}=n^{-1/3}\mathbf{S}$, où $n$ est le nombre de points et $\mathbf{S}$ la matrice de covariance de leurs coordonnées logarithmiques. Les points proches contribuent fortement à la somme, les points éloignés très peu. Un point entouré de nombreux voisins reçoit donc un poids plus élevé. Ces poids décrivent la concentration des mesures dans le diagramme ; ils ne représentent pas leurs incertitudes.

Le coefficient de corrélation de Pearson $r$ permet de quantifier l'intensité et le signe de la corrélation dans chaque relation de Kennicutt--Schmidt.

Les incertitudes sur les paramètres ajustés sont estimées à partir de 2000 échantillons bootstrap, obtenus en tirant avec remise les paires de mesures de gaz et de SFR. La relation est ajustée à chaque nouvel échantillon. Les 16e et 84e centiles des valeurs ajustées définissent les bornes de l'intervalle d'incertitude à 68\%. Ces intervalles décrivent l'incertitude d'échantillonnage, sans inclure les incertitudes systématiques sur la conversion CO--H$_2$ et la calibration du SFR. Lorsque la relation est presque verticale et que la pente est mal contrainte, nous indiquons l'angle $\theta$ de la droite par rapport à l'horizontale. Les droites en tirets de la figure~\ref{fig:individual_galaxies} représentent les ajustements. Le tableau~\ref{tab:individual_ks_summary} donne les paramètres et leurs incertitudes.


\begingroup
\raggedbottom
\setlength{\intextsep}{6pt}
\section[Relations de Kennicutt--Schmidt des différentes composantes gazeuses]{Relations de Kennicutt--Schmidt\\des différentes composantes gazeuses}
\label{sec:ks_individual_results}
Pour comprendre lequel des composantes \HI\ est la plus proche de la formation stellaire, nous étudions d'abord les relations de Kennicutt--Schmidt de chaque galaxie à la résolution commune à ses traceurs. Nous comparons ensuite les galaxies naines et spirales en regroupant leurs mesures à une résolution physique de 800 pc.\\\indent
Les panneaux inférieurs de la figure~\ref{fig:individual_galaxies} montrent les relations de l'\HI\ total, étroit et large, ainsi que celle du H$_2$ lorsque les données sont disponibles. DDO 53 et DDO 154 ne sont pas représentées : elles ne fournissent respectivement qu'une et six mesures espacées d'un faisceau à l'intérieur de la limite optique adoptée, à 400 pc.

Les relations du gaz atomique varient fortement d'une galaxie à l'autre (tableau~\ref{tab:individual_ks_summary}). Ho II présente de faibles corrélations pour les trois traceurs d'\HI. Dans IC 2574 et NGC 0925, l'\HI\ total est le mieux corrélé au SFR ; parmi les composantes séparées, la corrélation de l'\HI\ étroit est plus forte que celle de l'\HI\ large. Dans NGC 2403, NGC 2976 et NGC 4214, l'\HI\ total et large suivent mieux le SFR que l'\HI\ étroit. Dans NGC 7793 également, l'\HI\ total est mieux corrélé que chacune des composantes séparées. L'importance relative des corrélations étroite et large dépend donc de la galaxie. Ce résultat rejoint celui de \citet{Young2003}, qui ne trouvent pas de lien simple entre la présence d'\HI\ de faible dispersion et le SFR.

Dans les galaxies spirales, la relation de la composante étroite est faible pour NGC 0628, NGC 3184 et NGC 6946. NGC 3351 présente une corrélation négative pour l'\HI\ étroit et une faible corrélation positive pour la composante large. Dans NGC 2903 et NGC 4736, les deux composantes suivent une tendance positive, plus marquée pour l'\HI\ large. Dans NGC 5194, la composante large est aussi mieux corrélée que la composante étroite, mais l'\HI\ total donne la plus forte corrélation atomique. L'intensité, voire le signe, de la corrélation entre l'\HI\ étroit et le SFR varie donc entre ces spirales.

\citet{Fbigiel2008} trouvent eux aussi qu'un faible intervalle de densités de surface d'\HI\ peut correspondre à une gamme beaucoup plus large de densités de surface du SFR. Ils interprètent la saturation de l'\HI\ comme une signature de la transition vers le gaz moléculaire : le gaz supplémentaire se trouve principalement sous forme moléculaire. Dans cette interprétation, des régions de densités de surface d'\HI\ similaires peuvent présenter des quantités de gaz moléculaire et des SFR très différents.

Ces comparaisons individuelles utilisent des cartes de gaz et de SFR ramenées à la même résolution au sein de chaque galaxie. La taille physique du faisceau varie cependant entre galaxies. Dans la plupart des spirales, les relations moléculaires sont positives et mieux définies que les relations atomiques, avec des pentes comprises entre 0.86 et 1.20. Nous retrouvons ainsi la relation moléculaire presque linéaire déjà établie par \citet{Fbigiel2008}, en utilisant une régression orthogonale pondérée par la densité des points plutôt que la bissectrice OLS. La relation moléculaire sert ici de référence pour évaluer les liens de l'\HI\ étroit et large avec la formation stellaire. Parmi les galaxies de faible masse, NGC 2976 présente une forte corrélation moléculaire et une pente plus élevée, tandis que NGC 2403 montre une corrélation modérée, approximativement linéaire. Les corrélations moléculaires de NGC 0925 et NGC 4214 sont plus faibles et leurs pentes moins bien contraintes.

Pour comparer l'échantillon à une même échelle physique, nous lissons les cartes à 800 pc et regroupons les mesures espacées d'un faisceau. Nous examinons séparément les galaxies naines, les spirales et l'ensemble de l'échantillon. La figure~\ref{fig:pooled_dwarfs_ks} révèle une différence entre les deux populations. Dans l'échantillon de naines, l'\HI\ étroit est mieux corrélé au SFR ($r=0.51$) que l'\HI\ total ($r=0.45$) ou large ($r=0.23$). Dans les spirales, l'\HI\ étroit présente la plus faible corrélation ($r=0.11$), contre $r=0.37$ pour l'\HI\ large et $r=0.34$ pour l'\HI\ total. Pour l'ensemble des galaxies, les coefficients sont respectivement $r=0.18$, 0.34 et 0.35 pour l'\HI\ étroit, large et total.

Cette différence pourrait s'expliquer par le fait que la largeur des raies ne trace pas les mêmes conditions du gaz dans les deux populations. Dans les naines, l'émission étroite pourrait sélectionner un gaz plus froid et moins turbulent, associé à des conditions favorables à la formation stellaire. Dans les spirales, la turbulence peut élargir l'émission du gaz froid proche des régions de formation stellaire. De tels mouvements non thermiques contribuent à la largeur des raies du gaz atomique froid \citep{Saury2014}. L'énergie injectée par les jeunes étoiles peut entretenir la turbulence de l'\HI\ dans les disques actifs \citep{Tamburro2009}. Une partie du gaz froid peut ainsi être classée dans la composante large, ce qui pourrait expliquer sa plus forte corrélation avec le SFR dans les spirales.


\begin{table}[p]
\centering
\caption{\textbf{Ajustements résolus des relations gaz--SFR pour les galaxies individuelles et l'échantillon combiné.} $\theta$ est l'angle entre la droite ajustée et l'horizontale ; il décrit l'inclinaison de la relation, avec $N=\tan\theta$. Pour le gaz moléculaire, $N$ est l'exposant de la loi de puissance. Le coefficient $r$ est le coefficient de corrélation de Pearson.}
\label{tab:individual_ks_summary}
\par\vspace{8pt}
\begingroup

\footnotesize
\setlength{\tabcolsep}{3pt}
\renewcommand{\arraystretch}{1.12}
\textbf{(a) Gaz atomique}\\[5pt]
\begin{tabular*}{\textwidth}{@{\extracolsep{\fill}}lrrrrrrr@{}}
\toprule
 & & \multicolumn{2}{c}{\HI\ total} & \multicolumn{2}{c}{\HI\ étroit} & \multicolumn{2}{c}{\HI\ large} \\
Galaxie & $n$ & $\theta$ [$^\circ$] & $r$ & $\theta$ [$^\circ$] & $r$ & $\theta$ [$^\circ$] & $r$ \\
\midrule
Ho II & 117 & $83.52^{+2.04}_{-2.06}$ & $0.26$ & $82.71^{+3.07}_{-2.40}$ & $0.18$ & $82.65^{+2.87}_{-2.90}$ & $0.23$ \\
IC 2574 & 455 & $79.94^{+1.09}_{-0.77}$ & $0.46$ & $72.12^{+2.08}_{-1.46}$ & $0.41$ & $82.87^{+1.28}_{-1.06}$ & $0.28$ \\
NGC 0925 & 380 & $81.87^{+0.79}_{-0.61}$ & $0.56$ & $80.96^{+1.45}_{-1.52}$ & $0.42$ & $83.23^{+1.33}_{-1.27}$ & $0.32$ \\
NGC 2403 & 382 & $79.46^{+1.11}_{-0.99}$ & $0.61$ & $82.90^{+1.11}_{-1.09}$ & $0.36$ & $74.00^{+1.19}_{-1.44}$ & $0.60$ \\
NGC 2976 & 97/96/97 & $69.94^{+0.88}_{-1.54}$ & $0.88$ & $67.44^{+0.67}_{-6.25}$ & $0.46$ & $70.31^{+1.30}_{-1.17}$ & $0.88$ \\
NGC 4214 & 88 & $79.50^{+2.06}_{-2.06}$ & $0.60$ & $81.61^{+4.12}_{-4.65}$ & $0.42$ & $76.84^{+2.15}_{-2.08}$ & $0.57$ \\
NGC 7793 & 352 & $78.26^{+0.79}_{-0.73}$ & $0.78$ & $74.29^{+1.55}_{-1.47}$ & $0.49$ & $79.80^{+0.97}_{-1.03}$ & $0.58$ \\
NGC 0628 & 428 & $86.99^{+0.84}_{-0.79}$ & $0.26$ & $-87.59^{+0.92}_{-1.19}$ & $-0.02$ & $81.44^{+1.40}_{-1.12}$ & $0.31$ \\
NGC 2903 & 181/177/175 & $74.83^{+1.28}_{-1.54}$ & $0.78$ & $70.02^{+2.68}_{-2.28}$ & $0.57$ & $75.92^{+1.54}_{-1.75}$ & $0.67$ \\
NGC 3184 & 611 & $84.89^{+0.50}_{-0.41}$ & $0.43$ & $-89.85^{+0.74}_{-0.81}$ & $0.04$ & $81.87^{+0.83}_{-0.80}$ & $0.44$ \\
NGC 3351 & 314 & $87.35^{+0.95}_{-0.51}$ & $0.01$ & $-85.74^{+1.33}_{-1.38}$ & $-0.27$ & $83.23^{+1.91}_{-1.52}$ & $0.18$ \\
NGC 4736 & 84 & $70.30^{+1.59}_{-0.93}$ & $0.93$ & $71.52^{+2.29}_{-1.21}$ & $0.84$ & $68.60^{+1.52}_{-1.05}$ & $0.93$ \\
NGC 5194 & 288 & $84.21^{+1.40}_{-1.28}$ & $0.42$ & $83.80^{+1.82}_{-1.39}$ & $0.26$ & $83.80^{+1.40}_{-1.37}$ & $0.38$ \\
NGC 6946 & 216 & $87.36^{+0.72}_{-0.92}$ & $0.36$ & $-89.83^{+2.09}_{-2.76}$ & $0.07$ & $86.83^{+1.15}_{-1.60}$ & $0.34$ \\
\midrule
Naines & 769 & $83.21^{+0.53}_{-0.63}$ & $0.45$ & $80.26^{+0.74}_{-0.91}$ & $0.51$ & $85.23^{+0.80}_{-0.87}$ & $0.23$ \\
Spirales & 2119 & $84.77^{+0.47}_{-0.49}$ & $0.34$ & $89.83^{+0.51}_{-0.59}$ & $0.11$ & $78.49^{+0.70}_{-0.65}$ & $0.37$ \\
Ensemble & 2888 & $84.50^{+0.39}_{-0.39}$ & $0.35$ & $87.17^{+0.47}_{-0.49}$ & $0.18$ & $80.95^{+0.57}_{-0.52}$ & $0.34$ \\
\bottomrule
\end{tabular*}
\par\vspace{8pt}
\textbf{(b) Gaz moléculaire}\\[5pt]
\begin{tabular*}{0.76\textwidth}{@{\extracolsep{\fill}}lrrr@{}}
\toprule
Galaxie & $n$ & $N$ & $r$ \\
\midrule
NGC 0925 & 257 & $2.12^{+0.86}_{-0.72}$ & $0.25$ \\
NGC 2403 & 259 & $0.92^{+0.05}_{-0.12}$ & $0.58$ \\
NGC 2976 & 86 & $1.55^{+0.12}_{-0.14}$ & $0.85$ \\
NGC 4214 & 66 & $2.32^{+0.60}_{-0.95}$ & $0.40$ \\
NGC 0628 & 387 & $1.05^{+0.04}_{-0.04}$ & $0.73$ \\
NGC 2903 & 146 & $1.02^{+0.04}_{-0.06}$ & $0.84$ \\
NGC 3184 & 522 & $1.20^{+0.05}_{-0.07}$ & $0.64$ \\
NGC 3351 & 264 & $1.15^{+0.09}_{-0.10}$ & $0.71$ \\
NGC 4736 & 79 & $1.20^{+0.06}_{-0.06}$ & $0.91$ \\
NGC 5194 & 276 & $0.90^{+0.11}_{-0.01}$ & $0.72$ \\
NGC 6946 & 212 & $0.86^{+0.05}_{-0.05}$ & $0.84$ \\
\midrule
Naines & 507 & $1.15^{+0.04}_{-0.17}$ & $0.35$ \\
Spirales & 1918 & $1.07^{+0.02}_{-0.02}$ & $0.73$ \\
Ensemble & 2425 & $0.84^{+0.02}_{-0.02}$ & $0.60$ \\
\bottomrule
\end{tabular*}
\par\vspace{5pt}
\footnotesize\raggedright
$n$ est le nombre de mesures espacées d'un faisceau utilisées dans l'ajustement.
\par
\endgroup
\end{table}

\FloatBarrier
\clearpage

\vspace*{\fill}
\begin{figure}[H]
\centering
\includegraphics[width=\linewidth]{images/HO_II-cnm_wnm.pdf}
\par\vspace{6pt}
\includegraphics[width=\linewidth]{images/ks_individual/HO_II-sfr_to_gas.pdf}
\caption[Cartes de densité de surface et relations de Kennicutt--Schmidt des galaxies individuelles]{\textbf{Ho II.} En haut : densités de surface de l'\HI\ total, cinématiquement froid et cinématiquement chaud. En bas : relations de Kennicutt--Schmidt de l'\HI\ total, étroit et large, et du H$_2$ lorsque celui-ci est inclus. La distribution grisée représente tous les pixels positifs cartographiés. Les cercles bleus correspondent aux mesures de gaz et de SFR espacées d'au moins une largeur de faisceau. Les droites en tirets représentent les régressions orthogonales pondérées par la densité des points bleus, lorsqu'un ajustement est disponible. Pour les relations entre l'\HI\ total, étroit ou large et le SFR, $\theta$ est l'angle de la droite avec l'horizontale. Pour la relation H$_2$--SFR, nous indiquons la pente $N$, ou $\theta$ si la droite est trop proche de la verticale pour déterminer une incertitude finie sur la pente. Les intervalles correspondent aux intervalles bootstrap centraux à 68\% ; $r$ est le coefficient de Pearson et $n_{\mathrm{fit}}$ le nombre de mesures utilisées dans l'ajustement H$_2$--SFR. Les lignes gris clair indiquent des temps d'épuisement constants.}
\label{fig:individual_galaxies}
\label{fig:hi_maps_HO_II}
\label{fig:ks_individual_HO_II}
\end{figure}
\vfill
\clearpage

\vspace*{\fill}
\begin{figure}[H]
\ContinuedFloat
\captionsetup{labelformat=continuedgalaxies,labelsep=none,list=no}
\centering
\textbf{IC 2574}\par\vspace{4pt}
\includegraphics[width=\linewidth]{images/IC2574-cnm_wnm.pdf}
\par\vspace{6pt}
\includegraphics[width=\linewidth]{images/ks_individual/IC2574-sfr_to_gas.pdf}
\caption[]{}
\label{fig:hi_maps_IC2574}
\label{fig:ks_individual_IC2574}
\end{figure}
\vfill
\clearpage

\vspace*{\fill}
\begin{figure}[H]
\ContinuedFloat
\captionsetup{labelformat=continuedgalaxies,labelsep=none,list=no}
\centering
\textbf{NGC 0925}\par\vspace{4pt}
\includegraphics[width=\linewidth]{images/NGC0925-cnm_wnm.pdf}
\par\vspace{6pt}
\includegraphics[width=\linewidth]{images/ks_individual/NGC0925-sfr_to_gas.pdf}
\caption[]{}
\label{fig:hi_maps_NGC0925}
\label{fig:ks_individual_NGC0925}
\end{figure}
\vfill
\clearpage

\vspace*{\fill}
\begin{figure}[H]
\ContinuedFloat
\captionsetup{labelformat=continuedgalaxies,labelsep=none,list=no}
\centering
\textbf{NGC 2403}\par\vspace{4pt}
\includegraphics[width=\linewidth]{images/NGC2403-cnm_wnm.pdf}
\par\vspace{6pt}
\includegraphics[width=\linewidth]{images/ks_individual/NGC2403-sfr_to_gas.pdf}
\caption[]{}
\label{fig:hi_maps_NGC2403}
\label{fig:ks_individual_NGC2403}
\end{figure}
\vfill
\clearpage

\vspace*{\fill}
\begin{figure}[H]
\ContinuedFloat
\captionsetup{labelformat=continuedgalaxies,labelsep=none,list=no}
\centering
\textbf{NGC 2976}\par\vspace{4pt}
\includegraphics[width=\linewidth]{images/NGC2976-cnm_wnm.pdf}
\par\vspace{6pt}
\includegraphics[width=\linewidth]{images/ks_individual/NGC2976-sfr_to_gas.pdf}
\caption[]{}
\label{fig:hi_maps_NGC2976}
\label{fig:ks_individual_NGC2976}
\end{figure}
\vfill
\clearpage

\vspace*{\fill}
\begin{figure}[H]
\ContinuedFloat
\captionsetup{labelformat=continuedgalaxies,labelsep=none,list=no}
\centering
\textbf{NGC 4214}\par\vspace{4pt}
\includegraphics[width=\linewidth]{images/NGC4214-cnm_wnm.pdf}
\par\vspace{6pt}
\includegraphics[width=\linewidth]{images/ks_individual/NGC4214-sfr_to_gas.pdf}
\caption[]{}
\label{fig:hi_maps_NGC4214}
\label{fig:ks_individual_NGC4214}
\end{figure}
\vfill
\clearpage

\vspace*{\fill}
\begin{figure}[H]
\ContinuedFloat
\captionsetup{labelformat=continuedgalaxies,labelsep=none,list=no}
\centering
\textbf{NGC 7793}\par\vspace{4pt}
\includegraphics[width=\linewidth]{images/NGC7793-cnm_wnm.pdf}
\par\vspace{6pt}
\includegraphics[width=\linewidth]{images/ks_individual/NGC7793-sfr_to_gas.pdf}
\caption[]{}
\label{fig:hi_maps_NGC7793}
\label{fig:ks_individual_NGC7793}
\end{figure}
\vfill
\clearpage

\vspace*{\fill}
\begin{figure}[H]
\ContinuedFloat
\captionsetup{labelformat=continuedgalaxies,labelsep=none,list=no}
\centering
\textbf{NGC 0628}\par\vspace{4pt}
\includegraphics[width=\linewidth]{images/NGC0628-cnm_wnm.pdf}
\par\vspace{6pt}
\includegraphics[width=\linewidth]{images/ks_individual/NGC0628-sfr_to_gas.pdf}
\caption[]{}
\label{fig:hi_maps_NGC0628}
\label{fig:ks_individual_NGC0628}
\end{figure}
\vfill
\clearpage

\vspace*{\fill}
\begin{figure}[H]
\ContinuedFloat
\captionsetup{labelformat=continuedgalaxies,labelsep=none,list=no}
\centering
\textbf{NGC 2903}\par\vspace{4pt}
\includegraphics[width=\linewidth]{images/NGC2903-cnm_wnm.pdf}
\par\vspace{6pt}
\includegraphics[width=\linewidth]{images/ks_individual/NGC2903-sfr_to_gas.pdf}
\caption[]{}
\label{fig:hi_maps_NGC2903}
\label{fig:ks_individual_NGC2903}
\end{figure}
\vfill
\clearpage

\vspace*{\fill}
\begin{figure}[H]
\ContinuedFloat
\captionsetup{labelformat=continuedgalaxies,labelsep=none,list=no}
\centering
\textbf{NGC 3184}\par\vspace{4pt}
\includegraphics[width=\linewidth]{images/NGC3184-cnm_wnm.pdf}
\par\vspace{6pt}
\includegraphics[width=\linewidth]{images/ks_individual/NGC3184-sfr_to_gas.pdf}
\caption[]{}
\label{fig:hi_maps_NGC3184}
\label{fig:ks_individual_NGC3184}
\end{figure}
\vfill
\clearpage

\vspace*{\fill}
\begin{figure}[H]
\ContinuedFloat
\captionsetup{labelformat=continuedgalaxies,labelsep=none,list=no}
\centering
\textbf{NGC 3351}\par\vspace{4pt}
\includegraphics[width=\linewidth]{images/NGC3351-cnm_wnm.pdf}
\par\vspace{6pt}
\includegraphics[width=\linewidth]{images/ks_individual/NGC3351-sfr_to_gas.pdf}
\caption[]{}
\label{fig:hi_maps_NGC3351}
\label{fig:ks_individual_NGC3351}
\end{figure}
\vfill
\clearpage

\vspace*{\fill}
\begin{figure}[H]
\ContinuedFloat
\captionsetup{labelformat=continuedgalaxies,labelsep=none,list=no}
\centering
\textbf{NGC 4736}\par\vspace{4pt}
\includegraphics[width=\linewidth]{images/NGC4736-cnm_wnm.pdf}
\par\vspace{6pt}
\includegraphics[width=\linewidth]{images/ks_individual/NGC4736-sfr_to_gas.pdf}
\caption[]{}
\label{fig:hi_maps_NGC4736}
\label{fig:ks_individual_NGC4736}
\end{figure}
\vfill
\clearpage

\vspace*{\fill}
\begin{figure}[H]
\ContinuedFloat
\captionsetup{labelformat=continuedgalaxies,labelsep=none,list=no}
\centering
\textbf{NGC 5194}\par\vspace{4pt}
\includegraphics[width=\linewidth]{images/NGC5194-cnm_wnm.pdf}
\par\vspace{6pt}
\includegraphics[width=\linewidth]{images/ks_individual/NGC5194-sfr_to_gas.pdf}
\caption[]{}
\label{fig:hi_maps_NGC5194}
\label{fig:ks_individual_NGC5194}
\end{figure}
\vfill
\clearpage

\vspace*{\fill}
\begin{figure}[H]
\ContinuedFloat
\captionsetup{labelformat=continuedgalaxies,labelsep=none,list=no}
\centering
\textbf{NGC 6946}\par\vspace{4pt}
\includegraphics[width=\linewidth]{images/NGC6946-cnm_wnm.pdf}
\par\vspace{6pt}
\includegraphics[width=\linewidth]{images/ks_individual/NGC6946-sfr_to_gas.pdf}
\caption[]{}
\label{fig:hi_maps_NGC6946}
\label{fig:ks_individual_NGC6946}
\end{figure}
\vfill
\clearpage

\vspace*{\fill}
\begin{figure}[H]
\centering
\includegraphics[width=\linewidth]{images/ks_individual/pooled_groups_ks.pdf}
\caption[Relations de Kennicutt--Schmidt combinées des naines, des spirales et de l'ensemble de l'échantillon]{Relations de Kennicutt--Schmidt combinées des galaxies naines (en haut), des spirales (au milieu) et de l'ensemble de l'échantillon (en bas), à une résolution physique de 800 pc. De gauche à droite : \HI\ total, cinématiquement froid, cinématiquement chaud, puis H$_2$. Les symboles et les ajustements pondérés par la densité des points sont ceux de la figure~\ref{fig:individual_galaxies}. Les relations \HI--SFR indiquent $\theta$, et les relations H$_2$--SFR, $N$.}
\label{fig:pooled_dwarfs_ks}
\end{figure}
\vfill
\clearpage
\FloatBarrier

\endgroup


\begingroup
\raggedbottom
\setlength{\parskip}{6pt}
\section[Efficacité de formation stellaire]{Efficacité de formation stellaire}
\label{sec:sfe}

Les relations de Kennicutt--Schmidt décrivent les variations conjointes des densités de surface du gaz et du SFR. Suivant \citet{Leroy2008}, nous examinons comment l'efficacité apparente de formation stellaire diffère entre les réservoirs d'\HI\ étroit et large, et comment elle évolue dans le disque. L'efficacité de formation stellaire (SFE) est définie comme le SFR divisé par la masse de gaz.

Nous la calculons à partir des nouvelles cartes des composantes étroite et large, à la même résolution spatiale que les comparaisons individuelles de Kennicutt--Schmidt. Pour une région contenant $N_{\mathrm{reg}}$ pixels valides :
\begin{equation}
\label{eq:regional_sfe}
\mathrm{SFE}_p = 10^{-6}\,
\frac{\displaystyle\sum_{i=1}^{N_{\mathrm{reg}}}\Sigma_{\mathrm{SFR},i} A_i}
{\displaystyle\sum_{i=1}^{N_{\mathrm{reg}}}\Sigma_{p,i} A_i},
\end{equation}
\noindent où $A_i$ est l'aire du pixel et $p$ le traceur gazeux. Le facteur $10^{-6}$ convertit la densité de surface du SFR de kpc$^{-2}$ en pc$^{-2}$, de sorte que la SFE est exprimée en yr$^{-1}$. Pour comparer les traceurs sur une même aire, nous utilisons les valeurs des mêmes pixels dans les cartes de gaz et de SFR.

\begin{table}[tbp]
\centering
\caption{\textbf{M\'ediane de $\operatorname{\mathbf{log}}_{\mathbf{10}}\mathbf{SFE}$}, en yr$^{-1}$, calculée dans l'ouverture retenue de chaque galaxie avant le regroupement de l'échantillon. Les intervalles couvrent les 68\% centraux de 10\,000 rééchantillonnages bootstrap de galaxies entières, sans inclure les incertitudes de calibration et de décomposition. Les lignes d'\HI\ utilisent neuf naines et sept spirales. Les lignes incluant le gaz moléculaire utilisent les mêmes pixels couverts en CO dans quatre naines et sept spirales.}
\label{tab:sfe}
\begingroup

\renewcommand{\arraystretch}{1.3}
\begin{tabular}{lrrr}
\toprule
Gaz au dénominateur & Naines & Spirales & Ensemble \\
\midrule
\HI\ total & $-9.69^{+0.21}_{-0.10}$ & $-9.16^{+0.17}_{-0.13}$ & $-9.39^{+0.14}_{-0.29}$ \\
\HI\ étroit & $-9.29^{+0.21}_{-0.15}$ & $-8.66^{+0.05}_{-0.33}$ & $-9.08^{+0.19}_{-0.28}$ \\
\HI\ large & $-9.47^{+0.20}_{-0.06}$ & $-8.99^{+0.24}_{-0.00}$ & $-9.13^{+0.14}_{-0.27}$ \\
H$_2$ & $-8.45^{+0.19}_{-0.09}$ & $-9.10^{+0.13}_{-0.13}$ & $-8.96^{+0.39}_{-0.15}$ \\
H$_2$ + \HI\ étroit & $-9.02^{+0.13}_{-0.23}$ & $-9.25^{+0.02}_{-0.10}$ & $-9.23^{+0.17}_{-0.12}$ \\
H$_2$ + \HI\ large & $-8.99^{+0.04}_{-0.27}$ & $-9.35^{+0.09}_{-0.03}$ & $-9.26^{+0.25}_{-0.13}$ \\
H$_2$ + \HI\ total & $-9.22^{+0.06}_{-0.28}$ & $-9.43^{+0.01}_{-0.02}$ & $-9.42^{+0.20}_{-0.03}$ \\
\bottomrule
\end{tabular}
\endgroup
\end{table}

La figure~\ref{fig:adopted_sfe} montre la SFE en fonction du rayon pour les différentes composantes gazeuses. Chaque point représente la SFE médiane des galaxies contribuant à un intervalle radial de largeur $0.1r_{25}$. Dans les naines, le réservoir d'\HI\ étroit semble former les étoiles plus efficacement que le réservoir large dans l'ensemble du disque. Dans les spirales, la composante étroite présente une efficacité plus élevée dans le disque interne ($r<0.5r_{25}$), tandis que les deux efficacités se rapprochent à plus grand rayon. Dans les deux populations, la SFE apparente de l'\HI\ total et de chacune des composantes diminue vers l'extérieur, en accord avec la décroissance observée par \citet{Leroy2008} dans les régions dominées par l'\HI.

Dans les spirales, la SFE moléculaire varie relativement peu avec le rayon, alors que l'ajout de l'une ou l'autre composante atomique introduit une décroissance radiale et réduit les valeurs de SFE. Les profils H$_2$+\HI\ étroit et H$_2$+\HI\ large sont similaires.

Dans les régions les plus internes des spirales, ces deux combinaisons donnent des SFE proches de celle du H$_2$ seul, car le gaz des régions couvertes en CO y est principalement moléculaire. Ajouter une composante atomique modifie donc peu le résultat. À plus grand rayon, sa contribution augmente~: l'ajout d'\HI\ abaisse davantage la SFE et les profils combinés s'écartent du profil moléculaire, conformément à la diminution radiale de la fraction moléculaire observée par \citet{Leroy2008}. Dans les naines couvertes en CO, les deux combinaisons présentent déjà une SFE inférieure à celle du H$_2$ seul dans les régions centrales. Contrairement aux centres des spirales, ces régions contiennent une contribution atomique importante. Celle-ci augmentant vers l'extérieur, les deux profils combinés s'éloignent davantage du profil moléculaire.

Le tableau~\ref{tab:sfe} résume les SFE médianes des échantillons de naines et de spirales.

\begin{figure}[H]
\centering
\includegraphics[width=\linewidth]{images/ks_individual/adopted_sfe_profiles.pdf}
\caption[Efficacité de formation stellaire avec les nouvelles cartes cinématiques]{Profils radiaux de SFE des naines (en haut) et des spirales (en bas). À gauche : SFR observé divisé par l'\HI\ total, étroit ou large. À droite : H$_2$ et ses combinaisons avec les composantes atomiques ou l'\HI\ total, aux mêmes positions couvertes en CO. Chaque point représente la médiane des mesures des galaxies dans un intervalle de $0.1r_{25}$. Les bandes ombrées couvrent les 16e--84e centiles de la distribution entre galaxies pour l'\HI\ total (à gauche) et le H$_2$ (à droite) ; elles ne représentent pas l'incertitude sur la médiane.}
\label{fig:adopted_sfe}
\end{figure}
\FloatBarrier
\clearpage
\endgroup

